\documentclass[a4paper,10pt,twoside,twocolumn,openany,bg=full,nodeprecatedcode]{dndbook}

\usepackage[italian]{babel}
\usepackage[utf8]{inputenc}
\usepackage{array}
\usepackage{caption}
\usepackage{graphicx}

% The template's artwork lives in the vendored template directory.
\graphicspath{{../../../tools/latex/dnd/img/}}

\captionsetup[table]{labelformat=empty,font={sf,sc,bf,},skip=0pt}

\title{\includegraphics[width=.22\textwidth]{../../../assets/images/Free5e-Logo-Gold-Ring.png}\\[16pt]\color{titlered} Vanguard\\[4pt]\large Una classe di Free5e}
\author{Free5e\\[4pt]\small Edizione italiana}
\date{}

\begin{document}

\frontmatter
\maketitle
\tableofcontents
\mainmatter

\chapter{Vanguard}

\DndDropCapLine{I}{Vanguard incanalano la forza divina, ergendosi come campioni incrollabili di cause più grandi di loro.}

\begin{DndComment}[color=PhbLightCyan]{Violare il proprio giuramento}
  Un Vanguard cerca di mantenere i più elevati standard di condotta, ma anche il Vanguard più virtuoso è fallibile.

  Talvolta la via giusta si dimostra troppo impegnativa, talvolta una situazione richiede di scegliere il male minore e talvolta il calore delle emozioni induce un Vanguard a trasgredire il proprio giuramento.

  Un Vanguard che ha infranto un voto cerca in genere l'assoluzione da un Chierico che condivida la sua fede o da un altro Vanguard dello stesso ordine.
  Il Vanguard potrebbe trascorrere un'intera notte in veglia e preghiera come segno di penitenza, oppure intraprendere un digiuno o un atto simile di rinuncia.
  Dopo un rito di confessione e perdono, il Vanguard ricomincia da capo.

  Se un Vanguard viola volontariamente il proprio giuramento e non mostra alcun segno di pentimento, le conseguenze possono essere più gravi.
  A discrezione del Conduttore, un Vanguard impenitente potrebbe essere costretto ad abbandonare questa classe e ad adottarne un'altra.
\end{DndComment}

\section{Privilegi di classe}

In quanto Vanguard, ottieni i seguenti privilegi di classe.

\subsection{Punti ferita}

\textbf{Dadi Vita:} \texttt{1d10} per ogni livello da Vanguard

\textbf{Punti ferita al 1\degree livello:} \texttt{10 + il tuo modificatore di Costituzione}

\textbf{Punti ferita ai livelli superiori:} \texttt{1d10 (o 6) + il tuo modificatore di Costituzione} per ogni livello da Vanguard dopo il 1\degree.

\subsection{Competenze}

\textbf{Armature:} Tutte le armature, scudi

\textbf{Armi:} Armi semplici, armi da guerra

\textbf{Strumenti:} Nessuno

\textbf{Tiri salvezza:} Saggezza, Carisma

\textbf{Abilità:} Scegli due tra \emph{Atletica}, \emph{Intuizione}, \emph{Intimidire}, \emph{Medicina}, \emph{Persuasione} e \emph{Religione}.

\subsection{Equipaggiamento}

Inizi con il seguente equipaggiamento, oltre a quello conferito dal tuo background:

\begin{itemize}
  \item (a) un'arma da guerra e uno scudo oppure (b) due armi da guerra;
  \item (a) cinque giavellotti oppure (b) un'arma semplice da mischia qualsiasi;
  \item (a) uno zaino da sacerdote oppure (b) uno zaino da esploratore;
  \item una cotta di maglia e un simbolo sacro.
\end{itemize}

\section{Il Vanguard}

\begin{table*}[t]
\begin{DndTable}[header=Il Vanguard]{>{\centering\arraybackslash}p{.10\linewidth} >{\centering\arraybackslash}p{.18\linewidth} X}
  Livello & Bonus di competenza & Privilegi \\
  1\degree & +2 & Senso divino, Imposizione delle mani \\
  2\degree & +2 & Stile di combattimento, Lancio degli incantesimi, Punizione divina \\
  3\degree & +2 & Salute divina, Giuramento del dovere \\
  4\degree & +2 & Aumento dei punteggi di caratteristica \\
  5\degree & +3 & Attacco extra \\
  6\degree & +3 & Aura di protezione \\
  7\degree & +3 & Privilegio del Giuramento del dovere \\
  8\degree & +3 & Aumento dei punteggi di caratteristica \\
  9\degree & +4 & -- \\
  10\degree & +4 & Aura di coraggio \\
  11\degree & +4 & Punizione divina migliorata \\
  12\degree & +4 & Aumento dei punteggi di caratteristica \\
  13\degree & +5 & -- \\
  14\degree & +5 & Tocco purificatore \\
  15\degree & +5 & Privilegio del Giuramento del dovere \\
  16\degree & +5 & Aumento dei punteggi di caratteristica \\
  17\degree & +6 & -- \\
  18\degree & +6 & Miglioramenti delle aure \\
  19\degree & +6 & Aumento dei punteggi di caratteristica \\
  20\degree & +6 & Privilegio del Giuramento del dovere
\end{DndTable}
\end{table*}

\subsection{Senso divino (1\degree livello)}

La presenza di un male potente si manifesta ai tuoi sensi come un odore nauseabondo, mentre un bene potente risuona nelle tue orecchie come musica celestiale.
Come azione, puoi aprire la tua coscienza per percepire tali forze.
Fino alla fine del tuo prossimo turno, conosci la posizione di ogni celestiale, immondo o non morto entro 60 piedi da te che non si trovi dietro copertura totale.
Conosci il tipo (celestiale, immondo o non morto) di ogni essere di cui percepisci la presenza, ma non la sua identità (per esempio, il diavolo delle fosse Orobas).
Entro lo stesso raggio, percepisci anche la presenza di ogni luogo od oggetto che sia stato consacrato o profanato, come con l'incantesimo \emph{Santificare}.

Puoi usare questo privilegio un numero di volte pari a \texttt{1 + il tuo modificatore di Carisma}.
Quando completi un riposo lungo, recuperi tutti gli utilizzi spesi.

\subsection{Imposizione delle mani (1\degree livello)}

Il tuo tocco benedetto può curare le ferite.
Possiedi una riserva di potere curativo che si ricarica quando completi un riposo lungo.
Con questa riserva, puoi ripristinare un numero totale di punti ferita pari a \texttt{il tuo livello da Vanguard} \(\times 5\).

Come azione, puoi toccare una creatura e attingere potere dalla riserva per ripristinare a quella creatura un numero di punti ferita fino all'ammontare massimo ancora disponibile nella riserva.

In alternativa, puoi spendere 5 punti ferita dalla riserva di guarigione per curare il bersaglio da una malattia o neutralizzare un veleno che lo affligge.
Puoi curare più malattie e neutralizzare più veleni con un singolo utilizzo di Imposizione delle mani, spendendo punti ferita separatamente per ciascuno di essi.

Questo privilegio non ha effetto sui non morti e sui costrutti.

\subsection{Stile di combattimento (2\degree livello)}

Adotti uno stile di combattimento come tua specialità.
Scegli una delle opzioni seguenti.
Non puoi scegliere un'opzione dello Stile di combattimento più di una volta, anche se in seguito puoi scegliere di nuovo.

\subsubsection{Difesa}
Mentre indossi un'armatura, ottieni un bonus di +1 alla Classe Armatura.

\subsubsection{Duello}
Quando impugni con una mano un'arma da mischia senza impugnare altre armi, ottieni un bonus di +2 ai tiri per i danni con quell'arma.

\subsubsection{Combattere con armi possenti}
Quando tiri 1 o 2 su un dado dei danni per un attacco effettuato con un'arma da mischia che impugni con due mani, puoi ritirare il dado e devi usare il nuovo risultato.
L'arma deve avere la proprietà a due mani o versatile per ottenere questo beneficio.

\subsubsection{Protezione}
Quando una creatura che puoi individuare attacca un bersaglio diverso da te che si trova entro 5 piedi da te, puoi usare la tua reazione per imporre svantaggio al tiro per colpire.
Devi impugnare uno scudo.

\subsection{Lancio degli incantesimi (2\degree livello)}

Hai imparato ad attingere alla magia divina tramite la meditazione e la preghiera per lanciare incantesimi come fa un Chierico.

\textbf{Preparare e lanciare incantesimi}

La tabella seguente mostra quanti slot incantesimo possiedi per lanciare i tuoi incantesimi.

\begin{DndTable}[header=Slot incantesimo per mezzo incantatore]{>{\centering\arraybackslash}p{.18\linewidth} >{\centering\arraybackslash}X >{\centering\arraybackslash}X >{\centering\arraybackslash}X >{\centering\arraybackslash}X >{\centering\arraybackslash}X}
  Livello & 1\degree & 2\degree & 3\degree & 4\degree & 5\degree \\
  1\degree & -- & -- & -- & -- & -- \\
  2\degree & 2 & -- & -- & -- & -- \\
  3\degree & 3 & -- & -- & -- & -- \\
  4\degree & 3 & -- & -- & -- & -- \\
  5\degree & 4 & 2 & -- & -- & -- \\
  6\degree & 4 & 2 & -- & -- & -- \\
  7\degree & 4 & 3 & -- & -- & -- \\
  8\degree & 4 & 3 & -- & -- & -- \\
  9\degree & 4 & 3 & 2 & -- & -- \\
  10\degree & 4 & 3 & 2 & -- & -- \\
  11\degree & 4 & 3 & 3 & -- & -- \\
  12\degree & 4 & 3 & 3 & -- & -- \\
  13\degree & 4 & 3 & 3 & 1 & -- \\
  14\degree & 4 & 3 & 3 & 1 & -- \\
  15\degree & 4 & 3 & 3 & 2 & -- \\
  16\degree & 4 & 3 & 3 & 2 & -- \\
  17\degree & 4 & 3 & 3 & 3 & 1 \\
  18\degree & 4 & 3 & 3 & 3 & 1 \\
  19\degree & 4 & 3 & 3 & 3 & 2 \\
  20\degree & 4 & 3 & 3 & 3 & 2
\end{DndTable}

Per lanciare uno dei tuoi incantesimi da Vanguard di 1\degree livello o superiore, devi spendere uno slot del livello dell'incantesimo o superiore.
Recuperi tutti gli slot incantesimo spesi quando completi un riposo lungo.

Prepari la lista degli incantesimi da Vanguard che puoi lanciare, scegliendo dalla lista degli incantesimi da Vanguard riportata nei riferimenti di questo volume.
Quando lo fai, scegli un numero di incantesimi da Vanguard pari a \texttt{il tuo modificatore di Carisma + metà del tuo livello da Vanguard, arrotondata per difetto} (minimo un incantesimo).
Gli incantesimi devono essere di un livello per il quale possiedi slot incantesimo.

Per esempio, se sei un Vanguard di 5\degree livello, possiedi quattro slot incantesimo di 1\degree livello e due slot incantesimo di 2\degree livello.
Con un Carisma pari a 14, la tua lista degli incantesimi preparati può includere quattro incantesimi di 1\degree o 2\degree livello, in qualsiasi combinazione.
Se prepari l'incantesimo di 1\degree livello \emph{Cura ferite}, puoi lanciarlo usando uno slot di 1\degree o 2\degree livello.
Lanciare l'incantesimo non lo rimuove dalla tua lista degli incantesimi preparati.

Puoi cambiare la tua lista degli incantesimi preparati quando completi un riposo lungo.
Preparare una nuova lista di incantesimi da Vanguard richiede tempo trascorso in preghiera e meditazione: almeno 1 minuto per ogni livello dell'incantesimo, per ciascun incantesimo presente nella lista.

\textbf{Caratteristica da incantatore}

Il Carisma è la tua caratteristica da incantatore per i tuoi incantesimi da Vanguard, poiché il loro potere deriva dalla forza delle tue convinzioni.
Usi il tuo Carisma ogni volta che un incantesimo fa riferimento alla tua caratteristica da incantatore.
Inoltre, usi il tuo modificatore di Carisma quando determini la CD del tiro salvezza per un incantesimo da Vanguard che lanci e quando effettui un tiro per colpire con uno di essi.

\begin{DndComment}[color=PhbLightCyan]{Riferimenti rapidi}
  \textbf{CD del tiro salvezza dell'incantesimo} = \texttt{8 + il tuo bonus di competenza + il tuo modificatore di Carisma}

  \textbf{Modificatore di attacco dell'incantesimo} = \texttt{il tuo bonus di competenza + il tuo modificatore di Carisma}
\end{DndComment}

\textbf{Focus da incantatore}

Puoi usare un simbolo sacro come focus da incantatore per i tuoi incantesimi da Vanguard.

\subsection{Punizione divina (2\degree livello)}

Quando colpisci una creatura con un attacco con arma da mischia, puoi spendere uno slot incantesimo per infliggere danni radiosi al bersaglio, oltre ai danni dell'arma.
I danni extra sono \texttt{2d8} per uno slot incantesimo di 1\degree livello, più \texttt{1d8} per ogni livello dell'incantesimo superiore al 1\degree, fino a un massimo di \texttt{5d8}.
I danni aumentano di \texttt{1d8} se il bersaglio è un non morto o un immondo.

\subsection{Salute divina (3\degree livello)}

La magia divina che scorre dentro di te ti rende immune alle malattie.

\subsection{Giuramento del dovere (3\degree livello)}

Presti il giuramento che ti vincola per sempre come Vanguard.
Fino a questo momento ti sei trovato in una fase preparatoria, impegnato a seguire il cammino ma non ancora vincolato dal giuramento.
La tua scelta ti conferisce privilegi al 3\degree livello e di nuovo al 7\degree, 15\degree e 20\degree livello.
Questi privilegi includono gli incantesimi del giuramento e il privilegio Incanalare divinità.

\textbf{Incantesimi del giuramento}

Ogni giuramento possiede una lista di incantesimi associati.
Ottieni accesso a questi incantesimi ai livelli specificati nella descrizione del giuramento.
Una volta ottenuto accesso a un incantesimo del giuramento, lo hai sempre preparato.
Gli incantesimi del giuramento non contano ai fini del numero di incantesimi che puoi preparare ogni giorno.

Se ottieni un incantesimo del giuramento che non compare nella lista degli incantesimi da Vanguard, l'incantesimo è comunque un incantesimo da Vanguard per te.

\textbf{Incanalare divinità}

Il tuo giuramento ti permette di incanalare energia divina per alimentare effetti magici.
Ogni opzione di Incanalare divinità fornita dal tuo giuramento spiega come usarla.

Quando usi Incanalare divinità, scegli quale opzione usare.
Devi poi completare un riposo breve o lungo prima di poter usare di nuovo Incanalare divinità.

Alcuni effetti di Incanalare divinità richiedono tiri salvezza.
Quando usi un simile effetto di questa classe, la CD è pari alla CD del tiro salvezza dei tuoi incantesimi da Vanguard.

\subsection{Aumento dei punteggi di caratteristica (4\degree, 8\degree, 12\degree, 16\degree e 19\degree livello)}

Ai livelli 4, 8, 12, 16 e 19, puoi aumentare di 2 un punteggio di caratteristica oppure aumentare di 1 due punteggi di caratteristica, fino a un massimo di 20.

\subsection{Attacco extra (5\degree livello)}

Puoi attaccare due volte invece di una quando effettui l'azione Attacco durante il tuo turno.

\subsection{Aura di protezione (6\degree e 18\degree livello)}

A partire dal 6\degree livello, ogni volta che tu o una creatura amica entro 10 piedi da te dovete effettuare un tiro salvezza, la creatura ottiene un bonus al tiro salvezza pari al tuo modificatore di Carisma (con un bonus minimo di +1).
Devi essere cosciente per conferire questo bonus.
Al 18\degree livello, la gittata di questa aura aumenta a 30 piedi.

\subsection{Aura di coraggio (10\degree e 18\degree livello)}

A partire dal 10\degree livello, tu e le creature amiche entro 10 piedi da te non potete essere spaventati finché sei cosciente.

Al 18\degree livello, la gittata di questa aura aumenta a 30 piedi.

\subsection{Punizione divina migliorata (11\degree livello)}

Sei così permeato di potere virtuoso che tutti i tuoi colpi con armi da mischia portano con sé potere divino.
Ogni volta che colpisci una creatura con un'arma da mischia, la creatura subisce \texttt{1d8} danni radiosi extra.
Se usi anche Punizione divina con un attacco, aggiungi questi danni ai danni extra della tua Punizione divina.

\subsection{Tocco purificatore (14\degree livello)}

Puoi usare la tua azione per porre fine a un incantesimo su di te o su una creatura consenziente che tocchi.

Puoi usare questo privilegio un numero di volte pari al tuo modificatore di Carisma (minimo una volta).
Recuperi gli utilizzi spesi quando completi un riposo lungo.

\section{Giuramenti del dovere}

Diventare un Vanguard implica prestare voti che vincolano il Vanguard alla causa della rettitudine, un cammino attivo di lotta contro la malvagità.
Il giuramento finale, prestato quando raggiunge il 3\degree livello, è il culmine dell'addestramento del Vanguard.
Alcuni personaggi di questa classe non si considerano veri Vanguard finché non hanno raggiunto il 3\degree livello e prestato questo giuramento.
Per altri, prestare effettivamente il giuramento è una formalità, un timbro ufficiale su ciò che è sempre stato vero nel cuore del Vanguard.

\subsection{Giuramento di devozione}

Il Giuramento di devozione vincola un Vanguard ai più elevati ideali di giustizia, virtù e ordine.
Talvolta chiamati cavalieri, cavalieri bianchi o guerrieri sacri, i Vanguard incarnano l'ideale del cavaliere dall'armatura scintillante e agiscono con onore nella ricerca della giustizia e del bene superiore.
Si impongono i più elevati standard di condotta e alcuni, nel bene o nel male, impongono gli stessi standard al resto del mondo.
Molti di coloro che prestano questo giuramento sono devoti a divinità della legge e del bene e usano i principi delle loro divinità come misura della propria devozione.
Considerano ideali gli angeli, i perfetti servitori del bene, e incorporano immagini di ali angeliche nei loro elmi o stemmi.

\subsubsection{Principi della devozione (3\degree livello)}

Sebbene le parole esatte e i precetti del Giuramento di devozione varino, i Vanguard che prestano questo giuramento condividono i seguenti principi.

\begin{itemize}
  \item \textbf{Onestà.} Non mentire né barare. Fa' che la tua parola sia la tua promessa.
  \item \textbf{Coraggio.} Non temere mai di agire, sebbene sia saggio essere prudenti.
  \item \textbf{Compassione.} Aiuta gli altri, proteggi i deboli e punisci chi li minaccia. Mostra misericordia ai tuoi nemici, ma temperala con saggezza.
  \item \textbf{Onore.} Tratta gli altri con equità e lascia che le tue gesta onorevoli siano un esempio per loro. Compi quanto più bene possibile causando meno danni possibile.
  \item \textbf{Dovere.} Sii responsabile delle tue azioni e delle loro conseguenze, proteggi coloro che sono affidati alle tue cure e obbedisci a chi ha giusta autorità su di te.
\end{itemize}

\subsubsection{Incantesimi del giuramento (3\degree livello)}

Ottieni gli incantesimi del giuramento ai livelli da Vanguard indicati.

\begin{DndTable}[header=Incantesimi del Giuramento di devozione]{>{\centering\arraybackslash}p{.18\linewidth} X}
  Livello & Incantesimi da Vanguard \\
  3\degree & Protezione dal bene e dal male, Santuario \\
  5\degree & Ristorare inferiore, Zona di verità \\
  9\degree & Faro di speranza, Dissolvi magie \\
  13\degree & Libertà di movimento, Guardiano della fede \\
  17\degree & Comunione, Colonna di fuoco
\end{DndTable}

\subsubsection{Incanalare divinità (3\degree livello)}

Ottieni le due opzioni seguenti per Incanalare divinità.

\textbf{Arma sacra.} Come azione, puoi infondere energia positiva in un'arma che impugni usando Incanalare divinità.
Per 1 minuto, aggiungi il tuo modificatore di Carisma ai tiri per colpire effettuati con quell'arma (con un bonus minimo di +1).
L'arma emette inoltre luce intensa in un raggio di 20 piedi e luce fioca per altri 20 piedi oltre tale raggio.
Se l'arma non è già magica, diventa magica per la durata.

Puoi terminare questo effetto durante il tuo turno come parte di qualsiasi altra azione.
Se non impugni o non trasporti più quest'arma, oppure se diventi privo di sensi, l'effetto termina.

\textbf{Scacciare i profani.} Come azione, presenti il tuo simbolo sacro e pronunci o esegui con i segni una preghiera di condanna contro immondi e non morti, usando Incanalare divinità.
Ogni immondo o non morto che può percepirti entro 30 piedi da te deve effettuare un tiro salvezza su Saggezza.
Se la creatura fallisce il tiro salvezza, viene scacciata per 1 minuto o finché non subisce danni.

Una creatura scacciata deve usare i suoi turni per allontanarsi il più possibile da te e non può muoversi volontariamente in uno spazio entro 30 piedi da te.
Inoltre, non può effettuare reazioni.
Per la sua azione, può usare soltanto l'azione Scatto o tentare di liberarsi da un effetto che le impedisce di muoversi.
Se non c'è alcun luogo verso cui muoversi, la creatura può usare l'azione Schivata.

\subsubsection{Aura di devozione (7\degree e 18\degree livello)}

A partire dal 7\degree livello, tu e le creature amiche entro 10 piedi da te non potete essere affascinati finché sei cosciente.
Al 18\degree livello, la gittata di questa aura aumenta a 30 piedi.

\subsubsection{Purezza dello spirito (15\degree livello)}

Sei sempre sotto gli effetti di un incantesimo \emph{Protezione dal bene e dal male}.

\subsubsection{Nimbo sacro (20\degree livello)}

Come azione, puoi emanare un'aura di luce solare.
Per 1 minuto, da te si irradia luce intensa in un raggio di 30 piedi e luce fioca per altri 30 piedi oltre tale raggio.

Ogni volta che una creatura nemica inizia il suo turno nella luce intensa, subisce 10 danni radiosi.
Inoltre, per la durata, hai vantaggio ai tiri salvezza contro gli incantesimi lanciati da immondi o non morti.

Una volta usato questo privilegio, non puoi usarlo di nuovo finché non completi un riposo lungo.

\subsection{Giuramento dell'esemplare}

Il Giuramento dell'esemplare viene prestato da coloro che hanno incontrato il caos del mondo e scelto di prendere posizione.
Invece di urlare sopra il clamore, guidano con l'esempio e con l'azione, dimostrando che il mondo non ha bisogno di retorica infuocata o della manipolazione di parole melliflue, ma di persone che agiscano come dovrebbero, con determinazione e orgoglio.

I Vanguard che prestano questo giuramento possono provenire da molti background e sostenere infinite varianti della propria convinzione, ma più spesso tendono a essere \emph{Legali}.
Spesso ritengono che il potere delle parole risieda nella ragione e nella virtù che trasmettono, non nel volume o, peggio ancora, in magie manipolatrici che distorcono le menti di coloro su cui ricadono.

\subsubsection{Principi dell'esemplarismo (3\degree livello)}

Sebbene le parole esatte e i precetti del Giuramento dell'esemplare varino, i Vanguard che prestano questo giuramento condividono i seguenti principi.

\begin{itemize}
  \item \textbf{Temperanza.} Sussurrare nel silenzio trasmette più potere che urlare nel clamore.
  \item \textbf{Precisione.} Parla con determinazione e dì ciò che deve essere detto.
  \item \textbf{Santuario.} Sii il porto nella tempesta, sii la tranquillità nel caos.
  \item \textbf{Scopo.} Le azioni parlano più delle parole. Lascia che la tua spada e il tuo scudo siano la tua voce.
  \item \textbf{Sacro.} Diffida di coloro che usano la propria voce per piegare il mondo alla propria volontà per il proprio tornaconto e rimprovera coloro che usano la propria voce per piegare le creature alla propria volontà per il proprio tornaconto.
\end{itemize}

\subsubsection{Incantesimi del giuramento (3\degree livello)}

Ottieni gli incantesimi del giuramento ai livelli da Vanguard indicati.

\begin{DndTable}[header=Incantesimi del Giuramento dell'esemplare]{>{\centering\arraybackslash}p{.18\linewidth} X}
  Livello & Incantesimi da Vanguard \\
  3\degree & Santuario, Sonno \\
  5\degree & Sopprimere un senso, Silenzio \\
  9\degree & Controincantesimo, Dissolvi magie \\
  13\degree & Esilio, Divinazione \\
  17\degree & Santificare, Fuorviare
\end{DndTable}

\subsubsection{Incanalare divinità (3\degree livello)}

Ottieni le due opzioni seguenti per Incanalare divinità.

\textbf{Manto del silenzio.} Come azione bonus, presenti il tuo simbolo sacro e sussurri una preghiera, avvolgendo una creatura entro 60 piedi in un'aura di silenzio.
Per 10 minuti, il bersaglio non produce alcun rumore.
Non può parlare (compresi i componenti verbali degli incantesimi) e ottiene vantaggio alle prove di caratteristica di Destrezza (\emph{Furtività}).
Una creatura non consenziente effettua un tiro salvezza su Saggezza contro la CD del tiro salvezza dei tuoi incantesimi per impedire l'effetto quando viene bersagliata e può ripetere il tiro salvezza alla fine di ogni suo turno per porre fine all'effetto se lo supera.

\textbf{Santuario del sacrosanto.} Come azione, presenti il tuo simbolo sacro e sussurri una preghiera per offrire rifugio a chi fugge dal tumulto del mondo.
Tutte le creature a tua scelta entro 30 piedi ottengono punti ferita temporanei pari a \texttt{il tuo livello da Vanguard + il tuo modificatore di Carisma} e, se sono soggette all'effetto di una condizione o di un incantesimo contro il quale normalmente possono effettuare un tiro salvezza alla fine del loro turno, possono effettuare immediatamente un tiro salvezza contro l'effetto, ponendovi fine se lo superano, come di consueto.

\subsubsection{Aura di serenità (7\degree e 18\degree livello)}

A partire dal 7\degree livello, diffondi un'aura di quieta serenità che riduce il volume di tutti i rumori entro 10 piedi.
Tu e le creature amiche nell'aura ottenete resistenza ai danni da tuono e l'armatura non impone svantaggio alle prove di caratteristica di Destrezza (\emph{Furtività}) alle creature nell'aura.
Al 18\degree livello, la gittata di questa aura aumenta a 30 piedi.

\subsubsection{Rimprovero (15\degree livello)}

Quando una creatura entro la tua Aura di serenità tenta di parlare, attaccare o lanciare un incantesimo, puoi rimproverarla come reazione, facendole subire danni radiosi pari a \texttt{il tuo modificatore di Carisma} e costringendola a effettuare un tiro salvezza su Saggezza.
Se fallisce il tiro salvezza, il suo tentativo di parlare fallisce e la creatura rimane in silenzio fino all'inizio del suo prossimo turno.

\subsubsection{Mondo ordinato (20\degree livello)}

Al 20\degree livello, come azione, puoi infondere tranquillità e ordine nel mondo intorno a te per 1 minuto.
La tua Aura di serenità raddoppia la sua estensione e puoi usare Rimprovero senza spendere la tua reazione un numero di volte pari al tuo modificatore di Carisma, recuperando tutti gli utilizzi all'inizio del tuo turno successivo.
Per la durata, tutte le creature a tua scelta hanno vantaggio ai tiri salvezza contro le condizioni affascinato e spaventato o contro gli effetti degli incantesimi.

\chapter{Riferimenti del Vanguard}

\section{Lista degli incantesimi da Vanguard}

\subsection{1\degree livello}
\begin{itemize}
  \item Benedizione (concentrazione)
  \item Comando
  \item Cura ferite
  \item Individuazione del bene e del male (concentrazione)
  \item Individuazione del magico (concentrazione)
  \item Individuazione dei veleni e delle malattie (concentrazione)
  \item Favore divino
  \item Eroismo (concentrazione)
  \item Protezione dal bene e dal male (concentrazione)
  \item Purificare cibo e bevande
  \item Scudo della fede (concentrazione)
\end{itemize}

\subsection{2\degree livello}
\begin{itemize}
  \item Aiuto
  \item Punizione marchiante (concentrazione)
  \item Trova cavalcatura
  \item Riposo pacifico
  \item Ristorare inferiore
  \item Individuare oggetto (concentrazione)
  \item Arma magica (concentrazione)
  \item Preghiera di guarigione
  \item Protezione dal veleno
  \item Legame protettivo
  \item Zona di verità
\end{itemize}

\subsection{3\degree livello}
\begin{itemize}
  \item Creare cibo e acqua
  \item Luce diurna
  \item Dissolvi magie
  \item Cerchio magico
  \item Rimuovere maledizione
  \item Rianimare
\end{itemize}

\subsection{4\degree livello}
\begin{itemize}
  \item Bandire (concentrazione)
  \item Interdizione alla morte
  \item Individuare creatura (concentrazione)
\end{itemize}

\subsection{5\degree livello}
\begin{itemize}
  \item Dissolvi il bene e il male (concentrazione)
  \item Imposizione
  \item Rianimare morti
\end{itemize}

\end{document}
